\documentclass[../../../include/open-logic-section]{subfiles}

\begin{document}

\olfileid{sth}{ord-arithmetic}{using-addition}
\olsection{استخدام الجمع الترتيبي}

بجمع الأعداد الترتيبية نستطيع حساب رتب طائفة من المجموعات حسابًا
صريحًا، على معنى~\olref[spine][rank]{defnsetrank}:

\begin{lem}\ollabel{rankcomputation}
إذا كان~$\setrank{A} = \alpha$ و$\setrank{B} = \beta$، فإن:
\begin{enumerate}
	\item\ollabel{exrankpow} $\setrank{\Pow{A}} = \alpha\ordplus 1$
	\item\ollabel{exrankpair} $\setrank{\{A, B\}} = \max(\alpha,
	\beta) \ordplus 1$
	\item\ollabel{exrankcup} $\setrank{A \cup B} = \max(\alpha,
	\beta)$
	\item\ollabel{exranktuple} $\setrank{\tuple{A,B}} = \max(\alpha,
	\beta) \ordplus 2$
	\item\ollabel{exranktimes} $\setrank{A \times B} \leq \max(\alpha,
	\beta) \ordplus 2$
	\item\ollabel{exrankunion} $\setrank{\bigcup A} = \alpha$ عندما
	تكون~$\alpha$ خالية أو حدية؛ و$\setrank{\bigcup A} = \gamma$ عندما
	$\alpha = \gamma\ordplus 1$
\end{enumerate}
\end{lem}

\begin{proof}
نكرر استعمال~\olref[spine][rank]{ranksupstrict} في أجزاء البرهان.

\emph{\olref{exrankpow}.} إذا كان~$x \subseteq A$، كان
$\setrank{x} \leq \setrank{A}$، ومن ثم
$\setrank{\Pow{A}} \leq \alpha \ordplus 1$. ولما كان
$A \in \Pow{A}$، تعين~$\setrank{\Pow{A}} = \alpha \ordplus 1$.

\emph{\olref{exrankpair}.} يتبع من
\olref[spine][rank]{ranksupstrict}.

\emph{\olref{exrankcup}.} يتبع من
\olref[spine][rank]{ranksupstrict}.

\emph{\olref{exranktuple}.} طبق~\olref{exrankpair} مرتين.

\emph{\olref{exranktimes}.} لما كان
$A \times B \subseteq \Pow{\Pow{A \cup B}}$، استعمل
\olref{exrankpow} مرتين، ثم~\olref{exrankcup}.

\emph{\olref{exrankunion}.} إذا كان~$\alpha = \gamma\ordplus 1$،
وجد~$c \in A$ بحيث~$\setrank{c} = \gamma$، ولم يكن لأي
!!{element} من~$A$ رتبة أعلى؛ فيلزم~$\setrank{\bigcup A} = \gamma$.
وأما إذا كانت~$\alpha$ عددًا ترتيبيًا حديًا، فلـ$A$ !!{element}s
رتبها تقترب كيفما شئنا من~$\alpha$ مع بقائها دونه؛ ولذلك يكون
لـ$\bigcup A$ أيضًا !!{element}s رتبها تقترب كيفما شئنا من
$\alpha$ مع بقائها دونه. ومن ثم~$\setrank{\bigcup A} = \alpha$.
\end{proof}
\noindent
ويبقى بيان سبب ورود \emph{المتباينة} في~\olref{exranktimes} تمرينًا.

\begin{prob}
أعط مجموعتين~$A$ و$B$ بحيث
$\setrank{A \times B} =
\max(\setrank{A}, \setrank{B})$. وأعط مجموعتين~$A$ و$B$ بحيث
$\setrank{A \times B} =
\max(\setrank{A}, \setrank{B}) \ordplus 2$. هل توجد رتب أخرى ممكنة؟
\end{prob}

ونستطيع الآن أن نبين تكافؤ معان عدة معقولة لكون العدد الترتيبي
«منتهيًا» أو «غير منته».

\begin{lem}\ollabel{ordinfinitycharacter}
لأي عدد ترتيبي~$\alpha$، تتكافأ العبارات الآتية:
\begin{enumerate}
	\item\ollabel{ord:notinomega} $\alpha\notin \omega$، أي إن
	$\alpha$ ليس عددًا طبيعيًا؛
	\item\ollabel{ord:omegaplus} $\omega \leq \alpha$؛
	\item\ollabel{ord:oneplus} $1 \ordplus \alpha = \alpha$؛
	%\alpha \approx \alpha \ordplus 1$
	\item\ollabel{ord:plusone} $\alpha \approx \alpha\ordplus 1$، أي
	إن~$\alpha$ و$\alpha\ordplus 1$ متساويا العدد؛
	\item\ollabel{ord:infinite} $\alpha$ غير منتهة بحسب ديديكند.
	\end{enumerate}
\end{lem}
\noindent
فهذه خمسة أوصاف، متكافئة برهانيًا، لما يعنيه أن يكون العدد الترتيبي
غير منته، وبالنفي منتهيًا.

\begin{proof}
\emph{\olref{ord:notinomega} $\Rightarrow$ \olref{ord:omegaplus}.}
هذا من التثليث.

\emph{\olref{ord:omegaplus} $\Rightarrow$ \olref{ord:oneplus}.}
ثبت~$\alpha \geq \omega$. وبالاستقراء المتناهي التجاوز يوجد أصغر
عدد ترتيبي~$\gamma$، ولعله~$0$، بحيث يوجد عدد ترتيبي حدي~$\beta$
يحقق~$\alpha = \beta \ordplus \gamma$. وعندئذ:
\[
	1 \ordplus \alpha =
	1 \ordplus (\beta \ordplus \gamma) =
	(1 \ordplus \beta) \ordplus \gamma =
	\supstrict_{\delta < \beta} (1 \ordplus \delta) \ordplus \gamma =
	\beta \ordplus \gamma =
	\alpha.
\]
\emph{\olref{ord:oneplus} $\Rightarrow$ \olref{ord:plusone}.} من
الواضح وجود !!a{bijection}~$f \colon (\alpha \disjointsum 1) \to (1
\disjointsum \alpha)$. وإذا كان~$1 \ordplus \alpha = \alpha$، وجد
تماثل~$g \colon (1 \disjointsum \alpha) \to \alpha$. فانظر في
$\comp{f}{g}$.

\emph{\olref{ord:plusone} $\Rightarrow$ \olref{ord:infinite}.} إذا
كان~$\alpha \approx \alpha \ordplus 1$، وجدت !!a{bijection}
$f \colon (\alpha \disjointsum 1) \to \alpha$. عرّف
$g(\gamma) = f(\gamma, 0)$ لكل~$\gamma < \alpha$؛ فتكون هذه
!!{injection} شاهدة على أن~$\alpha$ غير منتهية بحسب ديديكند، لأن
$f(0,1) \in \alpha \setminus \ran{g}$.

\emph{\olref{ord:infinite} $\Rightarrow$ \olref{ord:notinomega}.}
وهو عين~\olref[z][infinity-again]{naturalnumbersarentinfinite}.
\end{proof}

\end{document}
